\section{The fixed sealed-extension calculus}

The preceding section stated the theorem informally: a sealed cubical extension must pay
primitive public trace through binary comparison, but higher structural trace is computed by
cubical horn filling.  This section defines the fixed object for which that statement is
proved.  The definitions are intentionally austere.  They do not try to model every feature
of Cubical Agda, every possible module system, or every way a library can grow.  They define
a raw sealed-extension calculus, denoted \(C_{\mathrm{ext}}\), in which payload, trace,
opacity, support, and trace cost can be measured without depending on a particular surface
presentation.

The section has three jobs.  First, it explains what counts as a sealed layer and what is
visible after sealing.  Second, it defines the normalized public signatures on which all
cardinality statements are made.  Third, it isolates the special fully coupled regime used by
the recurrence theorem.  The depth-two theorem itself is not proved here: the horn-computing
argument belongs to Sections~4 and~5.  What this section supplies is the accounting language
in which that theorem can be stated.

\subsection{Notation and terminology}

A library state is written \(B\).  Its sealed layers are written \(L_0,L_1,\ldots\), and
\(S(L_i)\) denotes the atomic public signature exported by \(L_i\) after counting
normalization.  A new declaration is written
\[
  e=(\Gamma_{\mathrm{new}},\Gamma_{\mathrm{act}},\Gamma_{\mathrm{cmp}},\sigma).
\]
The component \(\Gamma_{\mathrm{new}}\) contains payload, \(\Gamma_{\mathrm{act}}\) contains
unary action trace, \(\Gamma_{\mathrm{cmp}}\) contains comparison and horn-boundary trace, and
\(\sigma\) selects what is exported through the opaque public interface.

The following vocabulary is used throughout the paper.

\begin{center}
\begin{tabular}{ll}
\hline
Term & Meaning in this paper \\
\hline
Sealed extension & A new opaque layer added to the exported foundation or public interface. \\
Payload & The genuinely new mathematical or operational structure supplied by the layer. \\
Trace & Public compatibility data required for the payload to interact with prior exports. \\
Primitive trace & Trace not synthesized by normalization or cubical structural computation. \\
Derived trace & Trace present in a presentation but computed from lower data or generic operations. \\
Counting normal form & The canonical flattened telescope on which public fields are counted. \\
Active basis site & One irreducible primitive class of an earlier normalized public signature. \\
Full coupling & The regime in which the new layer acts on every active basis site in scope. \\
Sparse footprint & The general regime in which only selected active sites are touched. \\
Adequacy boundary & The bridge from richer semantic or surface systems into this raw calculus. \\
\hline
\end{tabular}
\end{center}

Two separations should be kept in mind.  First, payload may be higher-dimensional.  A user
may add higher inductive constructors, path constructors, or other genuinely higher content.
Those are counted as payload when they are part of the new layer's mathematical content.
Only the compatibility data forced by sealing the new layer against old public interfaces is
called structural trace.  Second, exact realized obligation objects and minimal public
signatures are different.  A higher horn obligation may remain present as a contractible exact
factor while contributing no primitive field to the minimal opaque signature.

\subsection{Raw syntax and judgments}

The theorem target is the following raw extension calculus.  Its grammar is deliberately
small:
\[
  i,j,n \in \mathrm{LayerId},
  \qquad
  B ::= \cdot \mid B,L_i,
  \qquad
  e ::= (\Gamma_{\mathrm{new}},\Gamma_{\mathrm{act}},\Gamma_{\mathrm{cmp}},\sigma).
\]
Here \(B\) is a monotone state of sealed layers.  The new declaration \(e\) has four parts.
The payload telescope \(\Gamma_{\mathrm{new}}\) introduces the new symbols of the layer.  The
action telescope \(\Gamma_{\mathrm{act}}\) describes unary action on selected old interface
sites.  The comparison telescope \(\Gamma_{\mathrm{cmp}}\) describes binary comparison,
transport, naturality, and horn-boundary data.  The export selection \(\sigma\) records which
elaborated fields are visible after sealing.

The fixed calculus is specified by the following judgments:
\[
\begin{array}{rcll}
  B\;\mathrm{wf} &&& \text{well-formed states},\\
  B\vdash e:\mathrm{RawDecl} &&& \text{raw declarations},\\
  B\vdash e\;\mathrm{adm} &&& \text{admissible declarations},\\
  B\vdash \mathrm{Elab}(e):\mathrm{Cand} &&& \text{elaboration to candidates},\\
  B\vdash \mathrm{Seal}(e):\mathrm{Layer} &&& \text{opaque sealing},\\
  B\vdash N(e):\mathrm{NormSig} &&& \text{normalized public signatures},\\
  B\vdash \phi\;\mathrm{primitive} &&& \text{primitive public fields},\\
  B\vdash \phi\;\mathrm{derived} &&& \text{transparent or computed fields},\\
  B\vdash \mathrm{Supp}(\phi)=S &&& \text{historical support of a field}. 
\end{array}
\]
An admissible declaration is therefore not just a syntactic tuple.  Its payload, action
clauses, comparison clauses, horn-boundary clauses, and export selection must type-check;
the selected footprint must match the advertised coupling mode; and the exported interface
must normalize to an opaque public signature.  A declaration lies outside the theorem target
when it uses forms not generated by this grammar, relies on hidden derivations after sealing,
advertises full coupling while omitting required footprint clauses, or asks for a general
Cubical Agda elaboration theorem rather than the fixed raw bridge of Section~3.

\begin{definition}[Raw structural clause classes]
The raw calculus uses three theorem-facing structural clause classes.

A \emph{raw action clause} is a field of \(\Gamma_{\mathrm{act}}\) whose support mentions one
prior opaque interface site and specifies the new payload's unary action on that site.

A \emph{raw comparison clause} is a field of \(\Gamma_{\mathrm{cmp}}\) whose support mentions
two prior interface sites and compares the two induced unary composites.  These include the
transport and naturality squares forced by old public paths or equivalences.

A \emph{raw horn clause} is a field of \(\Gamma_{\mathrm{cmp}}\) whose type presents a higher
structural boundary together with the missing-face and filler package of Section~5.  Raw horn
clauses may occur in exact realized obligation objects.  After the adequacy bridge and the
computational replacement theorem, they are classified as derived trace exactly when their
witness is synthesized by the cubical Kan operations from the lower boundary.
\end{definition}

The last sentence is important.  The word ``derived'' is not meant to be a mere syntactic tag.
The raw syntax records where a horn-shaped structural field appears; Sections~4 and~5 prove
that the relevant fields are computed by explicit cubical operations and therefore disappear
from \(\mu\)-minimal public trace.

\subsection{States, extensions, and public signatures}

\begin{definition}[State]
A \emph{state} or \emph{library} is a monotone context \(B\) closed under derivability.  An
evolution step
\[
  B_n \rightsquigarrow B_{n+1}
\]
adjoins one sealed integration layer, denoted \(L_{n+1}\).
\end{definition}

A state is not the same thing as a source file or an informal collection of lemmas.  It is the
opaque public interface available to future sealed extensions.  Transparent internal
construction may be used to build a layer, but after sealing only the selected normalized
fields of that layer are visible.

\begin{definition}[Surface extensions]
Over a state \(B_n\), a surface extension declaration is a finite quadruple
\[
  e=(\Gamma_{\mathrm{new}},\Gamma_{\mathrm{act}},\Gamma_{\mathrm{cmp}},\sigma),
\]
where \(\Gamma_{\mathrm{new}}\) is a telescope of newly declared payload symbols,
\(\Gamma_{\mathrm{act}}\) is a finite family of unary action clauses against the active
interface, \(\Gamma_{\mathrm{cmp}}\) is a finite family of comparison, transport, or
horn-boundary clauses, and \(\sigma\) specifies which elaborated fields are exported through
the sealed public interface.  The admissible extensions studied in this paper are exactly
those declarations whose intended action is global on the designated active interface
footprint.
\end{definition}

This is a fixed raw extension calculus, not the full surface language of Cubical Agda.  The
mechanized bridge checks the declaration forms just listed and their
admissibility/classification judgments.  It does not claim to parse or elaborate arbitrary
Agda programs.

\begin{definition}[Elaboration, sealing, and normalized public signature]
Each admissible surface extension declaration \(e\) over \(B_n\) elaborates to a core object
\[
  \mathrm{Elab}(e)
\]
of the fixed cubical core \(C\).  The payload telescope elaborates to new core declarations.
The action and comparison clauses elaborate to fields witnessing interaction with the active
interface.  The opacity component \(\sigma\) selects the public telescope exported after
sealing.

The \emph{normalized public signature} of \(e\) is the counting normal form
\[
  N(e) := N(\mathrm{Elab}(e)).
\]
It is obtained by flattening the exported public telescope, quotienting transparent telescope
isomorphisms, and marking as derived every field uniformly synthesized by definitional
equality, reflexivity, transport along already exported equalities, or generic structural
operations of the ambient cubical calculus.  Primitive trace cost is computed only after this
normalization.
\end{definition}

\begin{definition}[Presentation equivalence and minimal opaque cost]
Two admissible surface extensions \(e\) and \(e'\) over the same state are
\emph{presentation-equivalent}, written
\[
  e \approx e',
\]
when their elaborated sealed extensions determine the same public typing judgment up to
definitional equality, canonical telescope isomorphism, and transparent internal
normalization.

Write \(N_{\mathrm{tr}}(e)\subseteq N(e)\) for the irreducible primitive trace schemas in the
normalized public signature of \(e\).  The \emph{minimal opaque cost} of \(e\) is
\[
  \mu(e) := \min\{\, |N_{\mathrm{tr}}(e')| \mid e' \approx e \,\}.
\]
In the mechanization, the admissible presentation moves are generated by explicit
presentation steps, including reassociation, record splitting and bundling,
currying/uncurrying, transport of fields, transparent alias insertion or deletion, and
duplicate-derived deletion.
\end{definition}

The invariant \(\mu\) is deliberately insensitive to packaging.  A trace square does not become
cheaper by being hidden inside a one-constructor record, and it does not become more
expensive by being presented as several equivalent telescope spines.  What \(\mu\) counts is
irreducible primitive public trace after canonical normalization.

\begin{definition}[Candidate]
A \emph{candidate} over a library state \(B\) is a pair
\[
  X=(X_{\mathrm{core}},G_{\mathrm{obl}}),
\]
where \(X_{\mathrm{core}}\) is the mathematical content proposed for addition and
\(G_{\mathrm{obl}}\) is the family of normalized atomic structural obligations induced when
\(X\) is sealed against the active interface of \(B\).  Every admissible surface extension
declaration \(e\) determines such a candidate by elaboration.
\end{definition}

\begin{definition}[Counting normal form and atomic schemas]
Every public interface in this paper is counted only up to canonical telescope isomorphism:
reassociation and flattening of iterated dependent sums, splitting or rebundling a
single-constructor record into nested records, and currying or uncurrying iterated function
arguments when these preserve the same sealed typing judgment.  Concretely, an exported
interface is read as a dependent record, equivalently an iterated \(\Sigma\)-type whose field
types may themselves begin with \(\Pi\)-telescopes.  Its \emph{counting normal form} is the
fully flattened telescope obtained by reassociating the \(\Sigma\)-spine and expanding such
record wrappers.

An \emph{atomic schema} is one field of this counting normal form after transparent
telescope isomorphisms have been quotiented.  Such a field is tagged \emph{derived} when its
inhabitant is uniformly synthesized from earlier fields by definitional equality, reflexivity,
transport along already exported equalities, or the generic structural operations of the
ambient calculus.  It is tagged \emph{primitive} when no such transparent synthesis is
available in the fixed calculus.  The cost and basis notions below count primitive atomic
schemas; derived schemas may remain present in exact signatures but contribute no
primitive public trace unit.
\end{definition}

\begin{definition}[Core payload size]
For a candidate \(X\), let \(P(X)\) be the finite set of atomic payload schemas appearing in
the counting normal form of the public interface contributed directly by
\(X_{\mathrm{core}}\) before any structural witnesses against earlier layers are added.  Its
cardinality
\[
  \kappa(X) := |P(X)|
\]
is the \emph{core payload size} of \(X\).
\end{definition}

\begin{definition}[Sealed extension]
A \emph{sealed extension} is a candidate whose structural obligations have been discharged
and whose resolved data are exported only through an opaque public interface.  The opacity
invariant is the rule
\[
  \text{future layers may depend on } S(L_i),\text{ but not on the derivation of }\mathrm{Seal}(e_i).
\]
Thus future extensions may use exported payload and trace fields, but they may not inspect
the internal derivations that produced them.
\end{definition}

Opacity is what makes trace visible.  If future layers could freely inspect every proof term and
normalization derivation used to build \(L_i\), then some apparent public trace could be
recovered by peeking behind the interface.  A sealed extension disallows that move: the
normalized public signature is the contract with later layers.

\begin{definition}[Semantic sealed layer and public interface]
At the semantic level the \(n\)th sealed layer may be read as a package
\[
  E_n := (B_n,K_n,T_n,\operatorname{interp}_n).
\]
Here \(B_n\) is the incoming library state, \(K_n\) is the payload telescope contributed by
the new content, \(T_n\) is the resolved structural trace telescope exported after sealing, and
\(\operatorname{interp}_n\) interprets both parts in the ambient cubical foundation.  The public
interface of the sealed layer is the tagged coproduct
\[
  I_n := K_n \sqcup T_n.
\]
For the fixed raw instance, \(K_n\) is the payload part \(S_{\mathrm{core}}(L_n)\) of the
normalized public signature and \(T_n\) is the trace part \(S_{\mathrm{tr}}(L_n)\).  User-supplied
higher constructors and other mathematical content belong to \(K_n\) even when they are
genuinely higher-dimensional; only the integration data induced by sealing the layer against
earlier opaque interfaces belongs to \(T_n\).
\end{definition}

\begin{definition}[Promoted interface package]
A \emph{promoted interface package} over \(B\) is a finite package of constructively
irreducible generators and interaction data that is exported as one new sealed layer.  Such a
package need not modify the reduction rules of the kernel.  What matters is that its public
generators become part of the active exported interface for later sealing steps.
\end{definition}

\begin{definition}[Integration latency]
The \emph{integration latency} of a candidate \(X\) over \(B\) is
\[
  \Delta(X\mid B) := |G_{\mathrm{obl}}|,
\]
the number of irreducible atomic structural obligations that must be resolved before \(X\)
can be exported as a sealed layer.
\end{definition}

\begin{remark}[Refactoring invariance of the counting notions]
Because \(\kappa\), \(\Delta\), and later historical support are computed on counting normal
forms, admissible refactorings such as bundling laws into one record, splitting one interface
into several record shells, exporting curried rather than uncurried arguments, or pushing
transparently derivable witnesses deeper into an internal term do not change the counts.
Only a change in the irreducible primitive fields of the sealed typing judgment changes the
metrics.
\end{remark}

For a sealed layer \(L_k\) obtained from a candidate \(X_k\), we write \(S(L_k)\) for the set
of atomic exported schemas in the counting normal form of its public interface and
\[
  \kappa_k := \kappa(X_k)
\]
for its core payload size.  The indexing is chosen so that \(\Delta_k\) denotes the sealing cost
of \(L_k\).

\begin{proposition}[Transparent user-level code lies outside the recurrence]
Let \(U\) be a family of transparent definitions, lemmas, or record packages over \(B\) whose
symbols are definitionally reducible using the existing operational semantics and are not
promoted to a new opaque exported layer.  Then \(U\) contributes zero integration latency and
does not alter the active interface seen by later sealing steps.
\end{proposition}

\begin{proof}
Such a development creates no new irreducible interaction site against the exported
interface.  All reductions remain inherited from the pre-existing library, and no new opaque
generators are added to the global interface telescope.  Hence \(U\) is ordinary growth inside
one state, not a state transition of the form \(B_n\rightsquigarrow B_{n+1}\).
\end{proof}

\begin{definition}[Foundational core extension]
A candidate \(X\) over \(B_n\) is a \emph{foundational core extension} if its payload package is
intended to act globally on the currently active interface: the whole payload telescope
\(P(X)\) is operationally or semantically incomplete until its behavior has been specified
against every basis site in the designated coupling footprint inside the active exported
interface.  Typical examples are extensions of a universe-decoding apparatus, global
modalities presented by explicit modal dependent-type syntax whose context action must
distribute across existing type formers and transports, and heavily promoted generic
interfaces whose exported action is library-wide.
\end{definition}

\begin{remark}
Not every sealed extension is foundational in the sense of the previous definition.  A new
datatype or API layer may be orthogonal to much of the existing library and therefore require
only sparse interaction data.  Such extensions may preserve canonicity without becoming
globally coupled to earlier layers.  The recurrence theorem is intentionally a worst-case
statement about foundational core extensions, not a universality claim about all sealed
library growth.
\end{remark}

This scope restriction is essential.  Ordinary user-level additions typically enlarge a sparse
DAG of modules, datatypes, or APIs and do not force primitive interaction with unrelated
earlier exports.  The maximal-coupling results below isolate the special regime in which the
new sealed layer is advertised as a global operation on the active interface.  In that regime,
coverage of the advertised footprint is an exhaustiveness requirement rather than a stylistic
bookkeeping convention.

\subsection{Historical support, primitive classes, and bases}

The trace cost of a new layer depends not on the raw number of old fields but on the number
of old primitive classes that remain visible after normalization.  To make this invariant under
presentation, we choose basis representatives only after quotienting by transparent
derivability.

\begin{definition}[Historical interface]
For \(k\geq 1\), the \emph{historical interface of depth \(k\)} at stage \(n\) is
\[
  I_n^{(k)} := \coprod_{j=0}^{k-1} S(L_{n-j}),
\]
the tagged coproduct of the atomic schemas exported by the most recent \(k\) sealed layers
in counting normal form.
\end{definition}

The tags matter: if two layers happen to export isomorphic-looking fields, they are still
different public sites unless the normalized sealed interface identifies them.

\begin{definition}[Primitive classes and basis families]
For each sealed layer \(L\), let
\[
  \operatorname{PrimClasses}(S(L))
\]
be the finite set of primitive normal-form classes of the normalized public signature of
\(L\).  Two primitive fields represent the same class exactly when counting normalization
identifies them by transparent generation, definitional equality, uniform polymorphic
instantiation, transport along already exported equalities, or the generic structural
operations quotiented in Definition~2.7.

A \emph{basis family} \(S_{\mathrm{bas}}(L)\) is a chosen transversal of
\(\operatorname{PrimClasses}(S(L))\), that is, one representative primitive field from each
primitive normal-form class.  For \(k\geq 1\), once a basis family has been chosen for each of
the layers \(L_n,\ldots,L_{n-k+1}\), the associated \emph{historical basis of depth \(k\)} at
stage \(n\) is
\[
  I_{n,\mathrm{bas}}^{(k)} := \coprod_{j=0}^{k-1} S_{\mathrm{bas}}(L_{n-j}).
\]
Its elements are the basis sites on which a globally acting extension is primitively specified.
\end{definition}

\begin{theorem}[Canonical primitive classes determine bases]
For every sealed layer \(L\) in counting normal form, basis families in the sense of
Definition~2.18 exist.  If \(S_{\mathrm{bas}}(L)\) and \(S'_{\mathrm{bas}}(L)\) are two such basis
families, then projection to primitive normal-form classes identifies a canonical bijection
\[
  S_{\mathrm{bas}}(L) \cong S'_{\mathrm{bas}}(L).
\]
Consequently the cardinality \(|S_{\mathrm{bas}}(L)|\) is exactly
\(|\operatorname{PrimClasses}(S(L))|\), depends only on the canonical normalized public
signature of \(L\), and for every depth \(k\) the historical basis cardinality
\(|I_{n,\mathrm{bas}}^{(k)}|\) is independent of the chosen transversals and stable under
presentation-equivalent replacement of the participating layers.
\end{theorem}

\begin{proof}
The normalized public signature of \(L\) is finite, so the finite quotient
\(\operatorname{PrimClasses}(S(L))\) has a representative field in each class.  Choosing one
representative per class gives a basis family.  If two basis families are chosen, each
representative projects to exactly one primitive class, and each primitive class has exactly
one representative in each chosen transversal.  Matching representatives with the same class
gives the displayed canonical bijection.

Presentation-equivalent layers have canonically isomorphic normalized signatures by the
normalization theorem of Section~3.  That isomorphism preserves primitive/derived tags and
transparent-generation classes.  It therefore induces the same identification of primitive
class sets and hence the same historical basis cardinality.  No step uses the false principle
that arbitrary finite minimal generating families have canonical bijections; the invariant is
the primitive normal-form class set itself.
\end{proof}

\begin{lemma}[Primitive basis invariants]
For every sealed layer \(L\) and every presentation-equivalent replacement \(L'\), the
following data are canonically invariant: the finite primitive class set
\(\operatorname{PrimClasses}(S(L))\), the cardinality of any basis family
\(S_{\mathrm{bas}}(L)\), and the historical support and arity of every primitive trace field.
Consequently \(I_{n,\mathrm{bas}}^{(k)}\) has a presentation-independent cardinality.
\end{lemma}

\begin{proof}
The canonical normalization theorem of Section~3 gives a canonical isomorphism between the
normalized public signatures of presentation-equivalent layers.  The normalization relation
preserves primitive/derived tags, transparent-generation classes, and occurrences of opaque
layer identifiers in field types.  Therefore it induces a bijection on primitive class sets and
preserves historical support sets.  Basis families are transversals of those class sets by
Definition~2.18, so their cardinalities are invariant by Theorem~2.19.  Taking tagged
coproducts over the last \(k\) layers preserves these canonical bijections.
\end{proof}

\subsection{Global action, basis coverage, and canonical minimality}

This subsection isolates the hypothesis used by the recurrence theorem.  It is not a theorem
that every sensible extension is fully coupled.  It is a theorem about extensions that advertise
a total global action on a designated active footprint.  Once such a footprint is fixed,
coverage and minimality turn it into a precise trace count.

\begin{definition}[Basis coverage and canonical basis minimality]
An admissible foundational core candidate sealed against \(I_n^{(k)}\) has \emph{basis
coverage} if, for any historical basis \(I_{n,\mathrm{bas}}^{(k)}\) supplied by
Definition~2.18, every basis site receives at least one primitive interaction datum from the
new extension.  It is in \emph{canonical basis-minimal form} if, again for any such historical
basis, after elaboration and normalization each basis site contributes at most one irreducible
public trace schema.  By Theorem~2.19, these properties do not depend on which basis family
is chosen.
\end{definition}

\begin{definition}[Active footprint and full coupling judgment]
For a state \(B_n\) and a depth \(d\), an \emph{active footprint} is a finite subset
\[
  \mathrm{Foot}^{(d)}_n \subseteq I_{n,\mathrm{bas}}^{(d)}.
\]
A declaration is \emph{sparse at depth \(d\)} when its footprint is a proper subset of the
historical basis.  It is \emph{fully coupled at depth \(d\)}, written
\[
  B_n \vdash e\;\mathrm{fully\mbox{-}coupled}(d),
\]
when its admissibility judgment includes primitive action and comparison slots, modulo
canonical derivability, for every site of
\[
  \mathrm{Foot}^{(d)}_{n,\mathrm{full}} := I_{n,\mathrm{bas}}^{(d)}.
\]
Thus full coupling is an explicit hypothesis about the selected footprint, while sparse
footprints remain available for ordinary local or orthogonal growth.
\end{definition}

\begin{remark}[Trace-counting convention]
The recurrence uses one trace schema per active basis site per sealed layer.  A trace schema
may be a dependent record packaging the behavior of the whole new payload telescope at that
site; the individual clauses inside that package are not counted separately after bundling and
counting normalization.  A different convention, for example one primitive trace schema per
payload generator per active basis site, would insert an additional payload factor into the
cardinality theorem below and would change the recurrence.  This paper consistently uses the
bundled per-site convention.
\end{remark}

\begin{definition}[Fully coupled foundation]
A foundation is \emph{fully coupled} if the admissible sealed extensions under study are
foundational core extensions satisfying the full-coupling judgment of Definition~2.22 at the
relevant depth.  Equivalently, the paper's fully coupled declarations are globally acting
payload packages whose intended coupling footprint is the whole active historical basis, not
orthogonal additions with sparse dependency patterns.  This definition fixes the scope of the
action as a hypothesis.  The coverage and minimality theorems below derive the exact
irreducible primitive interaction data required to realize that action under the bundled
per-site trace convention of Remark~2.23.
\end{definition}

The next theorem turns basis coverage into a metatheoretic necessity for this globally coupled
regime rather than a stylistic choice.

\begin{theorem}[Full-coupling coverage in the fixed native-totality record]
Let \(B_n\) be a library state in a fully coupled foundation, let \(X\) be an admissible
foundational core extension sealed against the historical interface \(I_n^{(k)}\), and let
\(I_{n,\mathrm{bas}}^{(k)}\) be any historical basis for that interface.  If the extended theory is
to pass the sealed library's global admissibility discipline, then every basis site of
\(I_{n,\mathrm{bas}}^{(k)}\) must receive at least one primitive interaction datum from \(X\).  For
native core extensions this is enforced by canonicity-preserving totality of the operational
semantics; for promoted interface packages it is enforced by operational exhaustiveness of
the exported global interface.
\end{theorem}

\begin{proof}
By Definition~2.24, the admissible candidates considered here are not arbitrary sealed
additions but foundational core extensions whose payload is intended to act on the whole
active interface.  The theorem is therefore about that global regime only; orthogonal
extensions of the kind noted in Remark~2.16 are outside its hypotheses.

First consider a native global computational extension.  In a Tarski-style universe
presentation, suppose the active interface exports a code universe \(U\) and a decoding map
\(\operatorname{El}:U\to\mathrm{Type}\).  If the new sealed layer extends the universe with a
fresh code but the extended decoding machinery omits one basis site of the active interface,
then by Definition~2.18 there remains a family of raw exported generators whose behavior is
not recovered from the remaining primitive clauses by definitional equality, polymorphic
instantiation, transport, or generic structural operations.  Closed terms mentioning that
uncovered family no longer reduce through the advertised global operator, so the extension
is not total on the active interface it claims to cover.  In the native case this appears as a
canonicity or stuck-neutral failure for the extended language.  Thus native admissibility
forces at least one clause for every basis site of \(I_{n,\mathrm{bas}}^{(k)}\).

Now consider a promoted interface package.  Such a package need not alter kernel reduction,
so standard canonicity is not the right diagnostic.  What matters instead is operational
exhaustiveness of the exported global interface.  If the package is advertised as a
library-wide action, its public typing judgment is a total dependent map out of the record
enumerating the active interface.  Omitting one basis site leaves an irreducible family of
cases uncovered, because that site's action is not derivable from the other clauses by the
transparent operations built into Definition~2.18.  Later sealing steps cannot therefore treat
the package as a closed global operation.  Thus promoted packages also require at least one
primitive witness for every basis site of \(I_{n,\mathrm{bas}}^{(k)}\).
\end{proof}

\begin{proposition}[Canonical normalized presentations are basis-minimal]
Let \(e\) be an admissible surface extension declaration over \(B_n\) whose candidate
\(X(e)\) is sealed against the historical interface \(I_n^{(k)}\), and let
\(I_{n,\mathrm{bas}}^{(k)}\) be any historical basis for that interface.  Then the normalized public
signature \(N(e)\) contains at most one irreducible public trace schema for each basis site of
\(I_{n,\mathrm{bas}}^{(k)}\).  Equivalently, redundant primitive clauses that differ only by
transparent polymorphic compression, definitional equality, transport, or generic structural
operations do not survive in the canonical normalized presentation.
\end{proposition}

\begin{proof}
Suppose two primitive clauses of \(e\) correspond to the same basis site.  By construction of
Definition~2.18, they induce the same public trace field up to definitional equality,
transparent transport, uniform polymorphic instantiation, or the generic structural
operations already quotiented by counting normalization.  If both survived as distinct
irreducible fields in \(N(e)\), then the sealed public signature would contain
presentation-equivalent duplicate trace data.  But the normalization theorem of Section~3
retains exactly the irreducible visible public fields after normalization and identifies a
canonical normalized presentation for each presentation-equivalence class.  Hence one of the
two clauses is redundant and disappears in canonical normalized form.  Therefore at most one
irreducible trace schema survives per basis site.
\end{proof}

\begin{theorem}[Active-basis cardinality under full-coupling coverage]
Let \(B_n\) be a library state, let \(e\) be an admissible foundational core declaration over
\(B_n\), and suppose
\[
  B_n \vdash e\;\mathrm{fully\mbox{-}coupled}(k)
\]
in the sense of Definition~2.22.  Let \(X(e)\) be its elaborated candidate and let
\(I_{n,\mathrm{bas}}^{(k)}\) be any historical basis for the active interface.  Under the bundled
per-site convention of Remark~2.23,
\[
  \Delta(X(e)\mid B_n)
  = |\mathrm{Foot}^{(k)}_{n,\mathrm{full}}|
  = |I_{n,\mathrm{bas}}^{(k)}|.
\]
In particular, sealing cost counts irreducible basis clauses rather than the raw number of
exported generators, and the right-hand side is independent of the chosen basis family by
Theorem~2.19.
\end{theorem}

\begin{proof}
The judgment \(B_n\vdash e\;\mathrm{fully\mbox{-}coupled}(k)\) says that the selected active
footprint is the full historical basis
\[
  \mathrm{Foot}^{(k)}_{n,\mathrm{full}} = I_{n,\mathrm{bas}}^{(k)}.
\]
Coverage is exactly Theorem~2.25.  Minimality is exactly Proposition~2.26 after passing from
surface declarations to their canonical normalized public signatures.  Since
\(\Delta(X(e)\mid B_n)\) counts irreducible primitive obligations in that normalized
presentation, every basis site contributes one and only one irreducible primitive trace field.
Therefore
\[
  \Delta(X(e)\mid B_n)
  = |\mathrm{Foot}^{(k)}_{n,\mathrm{full}}|
  = |I_{n,\mathrm{bas}}^{(k)}|.
\]
\end{proof}

\subsection{The full-coupling trace decomposition principle}

In the fully coupled foundational-core regime, the preceding theorem identifies sealing cost
with the size of the active basis interface.  The next theorem explains how those resolved
obligations reappear in the next layer's public export alongside the layer's own new
mathematical payload.

\begin{theorem}[Full-coupling trace decomposition principle]
Each fully coupled sealed integration layer exports, in canonical basis form, a disjoint union
of two kinds of public schema: the new payload contributed by its core content and the trace
schemas corresponding to its resolved basis obligations.  Equivalently, if \(L_k\) is the layer
obtained by sealing a fully coupled candidate of core payload size \(\kappa_k\) and structural
cost \(\Delta_k\), then at the level of normalized public signatures there is a canonical
telescope isomorphism
\[
  N(L_k) \cong N_{\mathrm{core}}(L_k) \sqcup N_{\mathrm{tr}}(L_k),
\]
and this restricts on primitive normal-form classes to
\[
  \operatorname{Prim}(N(L_k))
  \cong
  \operatorname{Prim}(N_{\mathrm{core}}(L_k))
  \sqcup
  \operatorname{Prim}(N_{\mathrm{tr}}(L_k)).
\]
Thus, for any chosen basis transversal, there is a canonical decomposition
\[
  S_{\mathrm{bas}}(L_k) \cong S_{\mathrm{core}}(L_k) \sqcup S_{\mathrm{tr}}(L_k)
\]
with
\[
  |S_{\mathrm{core}}(L_k)|=\kappa_k,
  \qquad
  |S_{\mathrm{tr}}(L_k)|=\Delta_k.
\]
Hence
\[
  |S_{\mathrm{bas}}(L_k)|=\kappa_k+\Delta_k.
\]
\end{theorem}

\begin{proof}
There are two logically distinct sources of fields in the canonical normalized presentation.

First, the candidate's core content \(X_{k,\mathrm{core}}\) contributes genuinely new public
generators: a new type former, modality, promoted interface operator, universe rule, or
other primitive schema that did not previously exist in the library.  By definition of core
payload size, the number of these newly exported payload schemas is \(\kappa_k\).  These
form \(S_{\mathrm{core}}(L_k)\).

Second, by Theorem~2.27, sealing against the active interface requires exactly one irreducible
primitive interaction datum for each basis site of the active interface.  Formally, a sealed
layer is exported as a dependent record, equivalently an iterated \(\Sigma\)-type, whose fields
consist of the new payload together with the structural cells needed to type that payload
against earlier exports.  Opacity hides the inhabitants of those fields, but it does not hide
the field types occurring in the exported signature.  Thus once such a datum is supplied and
the layer is sealed, it no longer survives as an inspectable derivation visible to future
extensions; it survives as a public trace field recording how the new payload acts on an older
basis site.  These schemas form \(S_{\mathrm{tr}}(L_k)\), and there is one of them for each
resolved basis obligation.  Hence \(|S_{\mathrm{tr}}(L_k)|=\Delta_k\).

The two kinds of field are disjoint in the normalized telescope.  Payload symbols are
primitive generators introduced by the new layer itself, whereas trace schemas are indexed by
the prior basis interface against which the new layer was sealed.  Therefore
\[
  N(L_k) \cong N_{\mathrm{core}}(L_k) \sqcup N_{\mathrm{tr}}(L_k).
\]
Primitive and derived tags are part of counting normal form, so the same isomorphism
restricts to primitive normal-form classes.  Taking a basis transversal of the resulting
primitive classes gives
\[
  S_{\mathrm{bas}}(L_k) \cong S_{\mathrm{core}}(L_k) \sqcup S_{\mathrm{tr}}(L_k),
\]
and the cardinality formula follows immediately.
\end{proof}

\begin{remark}
The trace principle counts kernel-level extensions and promoted interface packages
uniformly.  In both cases the public export has a payload part and a trace part.  The
difference is the kind of primitive datum involved: for native extensions the payload
contributes new reduction heads and the trace contributes defining clauses, whereas for
promoted packages the payload contributes new exported operators and the trace contributes
comparison or transport witnesses.
\end{remark}

The section has now separated the quantities that are later combined.  The exact
depth-two theorem concerns the structural obligation object associated to a candidate.  The
\(\mu\)-minimality theorem concerns which trace fields survive as primitive public schemas
after normalization.  The recurrence theorem uses the special fully coupled case in which the
active footprint is the whole recent historical basis.  These are different claims, and the
remainder of the paper proves them in that order.
